\chapter{The Defence of the Old Order: I}
% \addcontentsline{toc}{chapter}{The Defence of the Old Order: I}

\section*{1}
In the last chapter, I noted that the classical Marxist writers had been well aware of the importance of tradition in shaping the consciousness of the working class—in shaping it, of course, in ways which made much more difficult a 'radical rupture' with the established order. But it is also necessary to note that classical Marxism did not really make very much of this phenomenon and that, with the signal exception of Gramsci, it did not seriously try to theorize, or even to identify, the many different ways in which the shaping of consciousness contributed to the stabilization and legitimation of capitalism.

This neglect may be related to that of political theory, which was discussed in the Introduction; and as in the case of political theory, the neglect was further accentuated in subsequent years by the kind of Marxism which Stalinism produced and was able to impose. This Marxism had a very pronounced tendency to provide whatever answer to any given question was most convenient to the people in power, which did not encourage the serious probing of difficult and often uncomfortable questions.

As a result, Marxism as a theory of domination remained poorly worked out. It was presented with the very large question of why capitalism was able to maintain itself, despite the crises and contradictions by which it was beset; and it tended to return a series of answers which were manifestly inadequate. In particular, it relied on an explanation based upon the Marxist view of the state as an instrument of capitalist coercion and repression. But coercion and repression could not possibly, in the case of many if not most of these regimes, explain why they endured. Nor did a second main line of explanation serve the purpose, namely betrayal by 'reformist' labour leaders, since this left whole the question why the working class allowed itself so regularly and so blatantly to be betrayed.

The notions of state coercion and repression, and of betrayal, are not wrong; but they need to be integrated into a wider theory of domination, comprising both 'infra-structural' and 'super- structural' elements. These elements can mostly be found in various parts of the corpus of Marxist writing, but have never been properly integrated. Until they are, a Marxist theory of politics will remain seriously deficient. The present chapter does not of course purport to fill the gap, but only to indicate some of the major ways in which the established order achieves legitimation.\footnote{This discussion is based on and continues the analysis which is to be found in The State in Capitalist Society (London, 1969).}

Tradition is not a monolith. On the contrary, it always consists of a large and diverse accumulation of customary ways of thought and action. In other words, there is not in any society one tradition but many: some of them are more congruent with others, some less. Thus, to take an instance which has direct relevance in the present context, there is in most societies a tradition of dissent as well as a tradition of conformity; or several of each. Traditional ways are never uniformly conservative.

But from a Marxist point of view, this 'polymorphous' nature of tradition is not particularly helpful. For however many forms it may assume, none of them is likely to afford a very helpful path to Marxist thought and to the revolutionary project which it proclaims: the Marxist notion of a 'most radical rupture' with traditional ideas, however attenuated it may be, signifies a break with all forms of tradition, and must expect to encounter the latter not as friend but as foe.

In this respect, the perpetuation of capitalism for well over a century after Marx and Engels began to write about it has produced an ironical twist in the story. One of the most important features of capitalism on which they fastened was precisely the universal uprooting of all aspects of life for which it was responsible. It is with something approaching exultation that they wrote in the Communist Manifesto that

\begin{displayquote}
    constant revolutionizing of production, uninterrupted disturbance of all social conditions, everlasting uncertainty and agitation distinguish the bourgeois epoch from ail earlier ones. All fixed, fast-frozen relations, with their train of ancient and venerable prejudices and opinions, are swept away, all new-formed ones become antiquated before they can ossify. All that is solid melts into air, all that is holy is profaned, and man is at last compelled to face with sober senses, his real conditions of life, and his relations with his kind.\footnote{Revs., p. 70. Note also, in the same vein, Marx's famous description of England's impact upon India, which 'separates Hindustan, ruled by Britain, from all its ancient traditions, and from the whole of its past history', and which 'produced the greatest, and, to speak the truth, the only social revolution ever heard of in Asia' ('The British Rule of India' in K. Marx and F. Engels, The First Indian War of Independence (Moscow, 1959), pp. 16, 19).}
\end{displayquote}

The clear message is of an irresistible movement, sweeping everything before it, including before long the forces which had set it going. In actual fact, the uprooting of tradition was not quite as thorough as Marx and Engels had claimed it to be; nor in any case were its consequences nearly as dramatic as they suggested. Furthermore, capitalism soon came to create, and has not ceased to reinforce, its own traditions, which were fused with or superimposed upon what remained of the older ones—and a lot did.

It was only a few years after the Manifesto that Marx, in the aftermath of the defeats of 1848, wrote of tradition weighing like a mountain upon the minds of the living, and of men making their own history not as they choose but 'under the given and inherited circumstances with which they are directly confronted;\endnote{The Eighteenth Brumaire of Louis Bonaparte in Surveys from Exile op cit p. 146.} and this was to be a recurrent theme with him. Indeed, we find in Capital two very different views of the impact of capitalist production itself upon the working class.

In the first of these, 'capitalist production develops a working class, which, by education, tradition, habit, looks upon the conditions of that mode of production as self-evident laws of Nature ... the dull compulsion of economic relations complete the subjection of the labourer to the capitalist';\endnote{K. Marx, Capital, I, p. 737.} and this process is much enhanced by the very nature of the capitalist mode of production which, much more than its predecessors, veils and mystifies the exploitative nature of its 'relations of production' by making them appear as a matter of free, unfettered, and equal exchange. The point is crucially important in the Marxist view of capitalist society and of its politics, and must be taken a little further.

There is an extremely strong sense in Marx of the falsity of perception which is woven into the tissue of capitalist society, the disjunction between appearance and reality, form and substance. In Capital, Marx writes as follows:
\begin{displayquote}
    If, as the reader will have realized to his great dismay, the analysis of the actual intrinsic relations of the capitalist process of production is a very complicated matter and very extensive; if it is a work of science to resolve the merely external movement into the true intrinsic movement, it is self-evident that conceptions which arise about the laws of production in the minds of agents of capitalist production and circulation will diverge drastically from these real laws and will merely be the conscious expression of the visible movements.\endnote{K. Marx, Capital, III, p. 307. My italics.}
\end{displayquote}

Marx then goes on to say that 'the conceptions of the merchant, stockbroker, and banker, are necessarily quite distorted'.

As for manufacturers, 'competition likewise assumes a completely distorted role in their minds'.\endnote{Ibid., p. 308.} But this distortion also occurs in a generalized form in the production of commodities, and engenders what Marx, in a famous section of Capi tal, called 'the fetishism of commodities'; and this 'fetishism' obviously affects those who produce the commodities, namely the workers. A commodity, Marx said, was 'a mysterious thing',
\begin{displayquote}
    simply because in it the social character of men's labour appears to them as an objective character stamped upon the product of that labour; because the relation of the producers to the sum total of their own labour is presented to them as a social relation, existing not between themselves, but between the products of their labour ... it is a definite social relation between men, that assumes, in their eyes, the fantastic form of a relation between things.\endnote{Capital, I, p. 72.}
\end{displayquote}

In the Grundrisse, Marx had made the same point, in a more general form still, with reference to the circulation of commodities produced as exchange values: 'The social relations of individuals to one another as a power over the individuals which has become autonomous, whether conceived as a natural force, as chance or in whatever other form, is a necessary result of the fact that the point of departure is not the free social individual.'\endnote{K. Marx, Grundrisse, p. 197.}

It is thus the capitalist system of production itself which also produces mystification as to the real nature of its 'relations of production'. This mystification is then further enhanced by intellectuals of one sort or another; and this incidentally indicates the role which revolutionary intellectuals should play, namely in helping to demystify capitalist reality.

There is, however, another view of the impact of capitalism in Marx. For he also believed that the working class becomes ever more 'disciplined, united, organized by the very mechanism of  the process of capitalist production itself', and that this must ultimately lead to a situation where 'the expropriators are expropriated'.\endnote{Capital, I, p. 763.}

These statements do not really contradict each other: they simply reflect different and contradictory facets of a complex reality, in which the opposing forces of tradition and actuality on the one hand, and of change on the other, do constant battle for the consciousness of the working class. From that battle, neither of these forces can emerge totally victorious, or totally secure in such victories as they may achieve. Tradition can never be completely paralysing: but neither can it be rapidly overcome. The problem, for victorious revolutions, is to prevent tradition from corroding them and ultimately defeating them from within. To topple a regime is seldom easy and is often very difficult indeed. But it is nevertheless easier to do so and to proclaim a new social order than actually to bring one into being. This is the point at which Lenin and Mao Tse-tung meet—in a common awareness that the revolution each led was under threat from the most deeply ingrained traditions of thought and behaviour. Lenin's last years were overshadowed by that awareness, and illness and death cut short whatever attempt he might have made to do more about it. Mao Tse-tung was more fortunate; but to what extent is far from clear.

At any rate, the enduring and pervasive importance of traditional ways, in a multitude of areas, is not in doubt, even in circumstances ol very rapid economic, social, and political change. Some of these ways are so deeply woven into the texture of life that they can survive more or less indefinitely without visible means of support and under conditions of extreme adversity. Two obvious cases in point are persecuted religious creeds and suppressed national sentiments. But for the most part, traditions are sustained and mediated by a network of particular institutions, which are actively involved in the performance of a process of transmission—institutions such as the family, schools, the mass media, and so on. It might be added that no reasonable government would seek to obliterate all traditions: the notion is both absurd and abhorrent.

The churches were the first mass media in history and the message they made it their business (in more senses than one) to transmit has generally tended to preach acceptance and obedience rather than questioning and rebellion. This is too summary for what is a complex and tortuous history; but it undoubtedly holds as a generalization. At any rate, Marxism, from the earliest days, has had a consistent record of opposition to religion and the churches, and this has been richly reciprocated by the latter. It was in the 'Introduction' to his Critique of Hegel's 'Philosophy of Right' that Marx made the famous statement that religion is 'the opium of the people',\endnote{K. Marx, Critique of Hegel's 'Philosophy of Right', p. 131.} a phrase which does not begin to do justice to the many different grounds of Marx's and Engels's opposition to religion and to their understanding of it as an historical fact, but which serves well enough to indicate what is in effect their paramount objection to it, namely that it stands as an obstacle to a proper appreciation by the working class of what is really wrong with the world it inhabits, so that 'the abolition of religion as the illusory happiness of the people is a demand for their true happiness.'\endnote{Ibid., p. 131.} 'The critique of religion disillusions man so that he will think, act, and fashion his reality as a man who has lost his illusions and regained his reason, so that he will revolve about himself as his own true sun.'\endnote{Ibid., p. 132.} The formulation is very much an 'early Marx' one: but the general sentiment remains at the core of the Marxist view of religion; so essentially does the view expressed in the Communist Manifesto that 'the parson has ever gone hand in hand with the landlord.'\endnote{Revs., p. 89.} Marx and Engels were equally scathing about the early (and for that matter the later) forms of Christian Socialism they encountered. But earnest attempts have been made over the years, and particularly in more recent times to attempt a 'reconciliation' between Christianity and Marxism by way of some kind of syncretic humanism. The results of these endeavours seem to have been philosophically and politically exceedingly thin. The question has of course nothing to do with the attempts of Communist regimes, particularly in Eastern Europe, to bring to an end their Kulturkampf with the churches and to find some mode of understanding with them.

A substantial part of Gramsci's concern with popular culture and with the degree to which it was permeated by a conservative 'common sense' had to do with the hold of religion upon the people, and with the need for Marxists to diffuse their own alternative 'common sense' as part of the battle for 'hegemony'. Religion was by no means the only part of the 'superstructure' with which he was concerned; but it was this which, in relation to popular culture, preoccupied him most, which is hardly surprising in the Italian context of the first decades of the twentieth century.

On the other hand, the notion of religion as the main ideological line of defence—or of attack, which is here the same thing—of the conservative forces in class society clearly does not accord with the ever-greater secular nature of this kind of society. Whether the religious 'opium' was ever as potent as it was declared to be in the first half of the nineteenth century (and before) is a question that cannot be answered in general: the impact of religion was obviously much greater on some parts of 'the people' than on others, depending on many different factors. But in any case, religion in advanced capitalist countries must now be reckoned, with obvious exceptions, to be one of the less effective forces which shape working-class consciousness, and certainly less than a good many others.

It is not my purpose here to discuss these forces in any detail, but only to indicate how they are, or might be, viewed in a Marxist perspective. The qualification is required because of the dearth of sustained Marxist work in analysing and exposing the meanings and messages purveyed in the cultural output produced for mass consumption in, say, the thirty-odd years since the end of World War II—not to speak of the virtual absence of such work in the years preceding it.\footnote{For a remarkable example of such analysis (by a non-Marxist), see George Orwell's essay on the reactionary values of boys' magazines in Britain, 'Boys' Weeklies', in G. Orwell, Collected Essays, Journalism and Letters, vol. I (London, 1970). The essay was written in 1939.} Gramsci was exceptional in his stress on the importance that must be attached to every artifact of this production, however trivial; and in his insistence that the struggle for 'hegemony' must be waged at every level, and include every single level of the 'superstructure'.

The key text on this whole issue was provided by Marx and Engels in The German Ideology, which was written in 1845—6, but which was only published in 1932 (save for one irrelevant chapter published earlier). In a Section entitled 'Concerning the Production of Consciousness', they wrote that
\begin{displayquote}
    The ideas of the ruling class are in every epoch the ruling ideas: i.e., the class, which is the ruling material force of society, is at the same time its ruling intellectual force. The class which has the means of material production at its disposal, has control at the same time over the means of mental production, so that thereby, generally speaking, the ideas of those who lack the means of mental production are subject to it.\endnote{K. Marx and F. Engels, The German Ideology (London, 1965), p. 60.}
\end{displayquote}

As will be argued presently, these formulations now need to be amended in certain respects. But there is at least one respect in which the text remains remarkably fresh, and points to one of the dominant features of life in advanced capitalist societies, namely the fact that the largest part of what is produced in the cultural domain in these societies is produced by capitalism; and is also therefore quite naturally intended to help, one way and another, in the defence of capitalism.

The point may be put quite simply: whatever else the immense output of the mass media is intended to achieve, it is also intended to help prevent the development of class- consciousness in the working class, and to reduce as much as possible any hankering it might have for a radical alternative to capitalism. The ways in which this is attempted are endlessly different; and the degree of success achieved varies considerably from country to country and from one period to another—there are other influences at work. But the fact remains that 'the class which has the means of material production at its disposal' does have 'control at the same time of the means of mental production'; and that it does seek to use them for the weakening of opposition to the established order. Nor is the point much affected by the fact that the state in almost all capitalist countries 'owns' the radio and television—its purpose is identical, namely the weakening of opposition to the established order. There is absolutely nothing remarkable about all this: the only remarkable thing is that the reality of the matter should be so befogged; and that Marxists should not have done more to pierce the fog which surrounds what is after all a vital aspect of the battle in which they are engaged.

Nor is that 'battle for consciousness' only waged in terms of traditions and communications. It has other facets which come under neither of these rubrics. One of the most important of these is the work process itself, to which Marx devoted some of his most searing pages—see for instance his denunciation of a system of division of labour which 'seizes labour-power by its very roots' and 'converts the labourer into a crippled monstrosity, by forcing his detail dexterity at the expense of a world of productive capabilities and instincts'.\footnote{K. Marx, Capital, I, p. 360. Here too, however, there is another side to the story. For 'Modern Industry', Marx also writes, 'compels society, under penalty of death, to replace the detail-worker of today, crippled by life-long repetition of one and the same trivial operation, and thus reduced to the mere fragment of a man, by the fully developed individual, fit for a variety of labours, ready to face any change of production, and to whom the different social functions he performs, are but so many modes of giving free scope to his own natural and acquired powers' (ibid., p. 488).}

Division of labour is one major aspect of the work process. Divisions within the working class, related to the division of labour, are another aspect of it. These concern skills, function, pay, conditions, status; and they add further to the erosion of class solidarity. Also, capitalist 'relations of production' are strongly hierarchical, not to say authoritarian (though this is not only true of capitalist ones), and subordination at work forms an important and diffuse—but complex and contradictory—element of working-class culture, going far beyond the work process. Its existence as a daily fact breeds frustrations which seek compensation and release in many different ways, most of them by no means conducive to the development of class- consciousness.

One of these ways is undoubtedly sport, or rather spectator and commercialized sport, some forms of which have assumed a central place in working-class life. For instance, vast numbers of people in the countries of advanced capitalism turn out on Saturdays and Sundays to watch soccer being played. Most of them are members of the working class, and so, in social origin, are the players, trainers, and managers. There is no other form of public activity which is capable of attracting even a fraction of those who attend football matches week in week out. A very large number of the people concerned are deeply involved, intellectually and emotionally, in the game, the players and one or other club; and their involvement, with all that informs and surrounds it, constitutes a sport culture which is an important part of the general culture. With this or that variation (e.g.
baseball instead of soccer in the United States), this is a major modern phenomenon, which radio and television have greatly helped to foster. The sport culture of capitalist countries, like every other mass activity, is big business for the various industries associated with sport, from sports gear and the pools industry to advertising. This is a very strong reason for the encouragement of its development by business.

But whether by design or not, there is an important amount of what might be described as cultural fall-out from the sports industry and spectator involvement. The nature of this fall-out is not quite as obviously negative as Marxists are often tempted to assume. The subject lends itself to simplistic and prim- sounding attitudes, which are often over-compensated by hearty demagogic-populist ones. In fact, the sport culture deserves, from the point of view of the making and unmaking of class- consciousness, much more attention than it has received. The elaboration of a Marxist sociology of sport may not be the most urgent of theoretical tasks; but it is not the most negligible of tasks either.\footnote{This would involve work on the evident 'sexist' bias and influence of most mass spectacle sports, for instance soccer, baseball and rugby, and of their celebration of 'manly' qualities. 'So what?' is a question to which it would be useful to have a detailed answer.}

The easiest assumption to make is that working-class involvement in sport as spectacle in the context of capitalism (the role and organization of sport in Communist countries raise questions of a different order) is most likely further to discourage the development of class-consciousness. But this is a bit too simple. For it rests on the antecedent assumption that a deep interest in the fortunes of, say, Leeds United Football Club is incompatible with militant trade unionism and the pursuit of the class struggle. This does not seem a priori reasonable and is belied by much evidence to the contrary; and to murmur 'bread and circuses5 is no substitute for serious thinking upon the matter.

What may well be advanced in relation to sport and the sport culture in capitalist countries is that it is very strongly pervaded by commercialism and money values; that this is very generally accepted and unquestioned as a 'natural5 part of the life of sport; and that it is also likely to strengthen an acceptance of social life in general as being 'naturally' and inevitably pervaded by commercialism and money values. In this sense, it may well be that the sport culture does help to block off the perception of a mode of social existence which is not thus pervaded. But how important this is in the total production of culture in these societies is a matter for surmise.

The formulations from The German Ideology which were quoted earlier have a serious defect in relation to present-day conditions, and indeed to earlier ones as well; and this defect is also shared by the Gramscian concept of 'hegemony', or at least by some interpretations of it.

What is involved is an over-statement of the ideological predominance of the 'ruling class', or of the effectiveness of that predominance. As 1 have noted earlier, it is at least as true now as it was when the words were written that 'the class which has the means of material production at its disposal has control at the same time over the means of mental production.' But it is only partially true—and the variations are considerable from country to country—that 'thereby, generally speaking, the ideas of those who lack the means of mental production are subject to it.' The danger of this formulation, as of the notion of 'hegemony', is that it may lead to quite inadequate account being taken of the many-sided and permanent challenge which is directed at the ideological predominance of the 'ruling class', and of the fact that this challenge, notwithstanding all difficulties and disadvantages, produces a steady erosion of that predominance.

A useful historical illustration of this process may be found in the ideological battle conducted in the eighteenth century against the ancien regime. The changes which occurred in the ideological climate of the ancien regime from year to year were almost imperceptible; but the difference between 1715, the year of the death of Louis XIV, and, say, 1775, is very great indeed. There are many closely related reasons for this: but one of them is the ideological battle that was conducted by a host of individuals, great and small, against the hegemony represented by the ancien regime. It is they who form the human link between 'infra-structure' and 'superstructure.'

On an enormously larger canvas, an ideological battle against bourgeois hegemony has been proceeding for 150 years and more; and even if only the relatively short span of the last fifty years is taken into account, it is obvious that, notwithstanding the dreadful battering to which it has been subjected, not least from its own side, what may for brevity's sake be called the socialist idea has vastly grown in strength in this period, immeasurably so in terms of its extension to every part of the globe, and also in the extent and depth of its penetration in the countries of advanced capitalism. However much more slowly and tortuously than he could ever have anticipated, Marx's 'old mole' has continued to burrow—so much so that the real question is progressively coming to be what kind of socialism towards which it is burrowing and, as a related question, how it is to be realized.

At any rate, the discussion of hegemony and class- consciousness more than ever requires the inclusion of the concept of a battle being fought on many different fronts and on the basis of the tensions and contradictions which are present in the actual structures or work and of life in general in capitalism as a social formation. The manifestations of that battle are endlessly diverse;\footnote{This diversity lends added significance to Marx's distinction between 'the material transformations of the economic conditions of production' which can be 'determined with the precision of natural science', and 'the legal, political, religious, esthetic or philosophical—in short, ideological forms' in which men become conscious of conflict and fight it out. (K. Marx, Preface to A Contribution to the Critique of Political Economy, in SW 1968, p. 183.)} but the fact is that it does occur. The ideological terrain is by no means wholly occupied by 'the ideas of the ruling class': it is highly contested territory.

In capitalist societies with bourgeois democratic regimes, ideological struggles are mainly waged in and through institutions which are not part of the state system. This point has in recent years been controverted by Louis Althusser and others who have opposed to it Althusser's notion of 'Ideological State Apparatuses' (ISA), according to which a vast number of institutions involved in one way or another in the dissemination of ideology are not only 'ideological apparatuses' but 'state ideological apparatuses', which must be distinguished from the 'Repressive State Apparatus'. Althusser lists these 'Ideological State Apparatuses' as 'the religious ISA (the system of the different churches), the educational ISA (the system of the different public and private 'Schools'), the family ISA, the legal ISA, the political ISA (the political system, including the different Parties), the trade union ISA, the communications ISA (press, radio, and television, etc.), the cultural ISA (Literature, the Arts, sports, etc.)'.\footnote{L. Althusser, 'Ideology and Ideological State Apparatuses (Notes Towards an Investigation)', in Lenin and Philosophy (London, 1972) p. 143. The listing seems a little superfluous, since it is difficult to think of anything that has been left out.}

Calling all these 'state ideological apparatuses' is based on or at least produces a confusion between class power and state power, a distinction which it is important not to blur.

Class power is the general and pervasive power which a dominant class (assuming for the purpose of exposition that there is only one) exercises in order to maintain and defend its predominance in 'civil society'. This class power is exercised through many institutions and agencies. Some of these are primarily designed for the purpose, e.g. political parties of the dominant class, interest and pressure groups, etc. Others may not be specifically designed for the purpose, yet may serve it, e.g. churches, schools, the family. But whether designed for the purpose or not, they are the institutions and agencies through which the dominant class seeks to assure its 'hegemony'. This class power is generally challenged by a counter-power, that of the subordinate classes, often through the same institutions, and also through different ones. Instances of the former case are the family, schools and churches; and of the latter, parties of the working class and trade unions. That some institutions are 'used' by opposite classes simply means that these institutions are not 'monolithic' but that on the contrary they are themselves arenas of class conflict; and some institutions which are designed to further class power may well play another role as well (or even altogether), for instance that of 'routinizing', stabilizing, and limiting class conflict. But this too only shows that institutions may assume different roles than those for which they were brought into being, and may have other effects than those which they were intended to have.

The important point, however, is that the class power of the dominant class is not exercised, in some important respects, by state action but by class action, at least in 'bourgeois democratic' regimes, and in a number of other forms of capitalist regimes as well. Of course, there are other, crucial respects, in which that power is exercised through the agency of the state; and the state is in all respects the ultimate sanctioning agency of class power. But these are different issues. I am concerned to point here to the fact that the dominant class, under the protection of the state, has vast resources, immeasurably greater than the resources of the subordinate classes, to bring its own weight to bear on 'civil society'. The following quotation from Marx's Class Struggles in France provides a good illustration of the meaning of class power. Writing about the election campaign of 'the party of Order' in 1849, Marx notes:
\begin{displayquote}
    The party of Order had enormous financial resources at its disposal: it had organised branches throughout France; it had all the ideologists of the old society in its pay; it had the influence of the existing government power at its disposal; it possessed an army of unpaid vassals in the mass of the petty bourgeoisie and peasants, who, still separated from the revolutionary movement, found in the high dignitaries of property the natural representatives of their petty property and their petty prejudices. Represented throughout the country by innumerable petty monarchs, the party of Order could punish the rejection of its candidates as insurrection, and could dismiss rebellious workers, recalcitrant farm labourers, servants, clerks, railway officials, registrars and all the functionaries who are its social subordinates.\endnote{SE, p. 90.}
\end{displayquote}

Clearly, much of this would now have to be amended, extended or qualified, and much that is more 'modern' by way of resources would have to be added: but the basic conception remains valid and important about this kind of capitalist regime, as opposed to other kinds, where class power is mainly or mostly exercised by way of the state. The extreme example of class power in capitalist society being 'taken over' and exercised by the state is Fascism and Nazism, with other forms and degrees of authoritarian rule in between.

In periods of acute social crisis and conflict, class power does tend to be taken over by the state itself, and indeed gladly surrenders to it; and it may well be argued that, even in the 'normal' circumstances of advanced capitalism, the state takes over more and more of the functions hitherto performed by the dominant class, or at least takes a greater share in the performance of these functions than was the case in previous periods. I am not here referring to state intervention in the economy, where this is obviously true,\endnote{See below, pp. 92 ff.} but rather to the shaping of consciousness.

The enormous advances in communication techniques have in any case tempted governments to use these techniques in order to try and influence 'public opinion', and to intervene much more directly in the shaping of consciousness. This is done both positively and negatively. Positively in the sense that there now exists a governmental and state communications industry, which actually disseminates or influences the dissemination of news, views, opinions and perspectives; and negatively in so far as governments try as best they can to prevent the dissemination of news, views etc. which they consider 'unhelpful'. The range of means and techniques for doing this varies from country to country but is everywhere formidable.

This communications industry run by the state does not pro-  duce any kind of commodity: its range of products, unlike its range of means and techniques, is fairly narrow, in ideological terms, and the products quite naturally bear mostly a conformist and helpful label—helpful, that is, to the maintenance of stability, the discrediting of dangerous thoughts and the defence of things as they are. The state is now in the consciousness industry in a big way, and thus plays its role in a permanent battle of ideas which is a crucial element of class conflict.

Even so, and even if the process of 'statisation' is taken fully into account, as it obviously must be, there is absolutely no warrant for speaking of 'state ideological apparatuses' in regard to institutions which, in bourgeois democratic societies, are not part of the state; and much which is important about the life of these societies is lost in the obliteration of the distinction between ideological apparatuses which are mainly the product of civil society' and those which are the product and part of the state apparatus. The point, it may be added, does not only bear on capitalist societies: it has large implications for the way in which ideological apparatuses are viewed in relation to a future socialist society, their connection to and independence from the state, and so on. As Althusser might say, these are not 'innocent' discussions.

\section*{2}

If it is agreed that a crucially important battle for consciousness is waged, or occurs, day in day out in capitalist societies (and in others as well for that matter), it also follows that a great deal of imporiance must be attributed to those people who play the main part in articulating and giving expression to the terms of that battle—intellectuals.

Intellectual is a notoriously difficult designation. Gramsci noted that all men were in a way 'intellectuals'. 'In any physical work, even the most degraded and mechanical', he wrote, 'there exists a minimum of technical qualification, that is, a minimum of creative intellectual activity';\endnote{A. Gramsci, Selections from the Prison Notebooks (London 1971) p. 8.} and, rather more positively,
\begin{displayquote}
    there is no human activity from which every form of intellectual participation can be excluded: homo faber cannot be separated from homo sapiens. Each man, finally, outside his professional activity, carries on some form of intellectual activity, that is. he is a 'philosopher', an artist, a man of taste, he participates in a particular conception of the world, has a conscious line of moral conduct, and therefore contributes to sustain a conception of the world or to modify it, that is, to bring into being new modes of thought.\endnote{Ibid., p. 9.}
\end{displayquote}

This is broadly speaking true and points to the fact that everyone, in his and her daily life and utterances, does mediate social reality for other people and thereby helps, in however so small a wTay, to shape the character of that reality. But Gramsci also observed that though all men might in some sense be said to be intellectuals, 'not all men have in society the function of intellectuals'.\endnote{Ibid., p. 9.}

As for that function, he defined it as being two-fold, 'organisational and connective': 'The intellectuals are the dominant group's “deputies” exercising the subaltern function of social hegemony and political government'; in other words, as involving two major superstructural levels, one of these being the organization of 'hegemony' and the second being the organization of political functions.\endnote{Ibid., p. 12.}

'This way of posing the problem', Gramsci said, 'has as a result a considerable extension of the concept of intellectual, but it is the only way which enables one to reach a concrete approximation of reality.'\endnote{Ibid., p. 12.}

This is not particularly convincing; and it would seem rather less confusing, without any loss, to confine the notion of 'intellectual' to the first of Gramsci's two categories and levels, namely to the cultural-ideological domain.\footnote{There is another important distinction of Gramsci's, that between 'organic' and 'traditional' intellectuals, which does not seem to be very helpful either. Roughly speaking, 'organic' intellectuals are those which a 'new class' needs because they 'give it homogeneity and an awareness of its own function not only in the economic but also in the social and political fields'. 'Traditional' intellectuals are those which are linked to already established classes. (See, ibid., p. 5 ff.) Gramsci had in mind the bourgeoisie and the landed aristocracy respectively, but the terminology and for that matter the conceptualization becomes unhelpful when the bourgeoisie is counterposed to the working class.} This is how Marx and Engels saw it, in The German Ideology, when they spoke of a 'division of labour' in the ruling class,
\begin{displayquote}
    so that inside this class one part appears as the thinkers of the class (its active, conceptive ideologists, who make the perfecting of the illusion of the class about itself their chief source of livelihood), while the others' attitude to these ideas and illusions is more passive and receptive, because they are in reality the active members of this class and have less time to make up illusions and ideas about themselves.\endnote{K. Marx and F. Engels, The German Ideology, p. 61.}
\end{displayquote}

The general view expressed here is that intellectuals/ ideologists are mainly (though not exclusively, as will be seen in a moment) the 'deputies' of the ruling class, in Gramsci's formula. Marx's own resounding phrase in Capital about economists was that the latter had, after the heyday of political economy, turned themselves into the 'hired prize-fighters' of the bourgeoisie;\endnote{K. Marx, Capital, I, p. 15.} and this is obviously capable of application to intellectuals and ideologists other than economists. Gramsci also called intellectuals the 'managers of legitimation', which fixes their role in this perspective with adequate precision; for it is indeed the management, fostering, and consolidation of the legitimacy of the existing social order that most intellectuals have in one way or another been about. The impression has often been purveyed in this century that intellectuals were for the most part left-wing dissidents; and the term 'intellectual' was itself invented in France at the time of the Dreyfus Affair to designate (pejoratively) those journalists or men of letters who were Dreyfusards' and who subversively attacked sacred institutions such as the French Army and cast doubt on its honour.

The Marxist assessment is the more realistic one: the majority of those who could properly be called intellectuals in bourgeois society, not to speak of professional people of one sort or another, lawyers, architects, accountants, doctors, scientists, etc., have been the 'managers of legitimation' of their society.

Nor have they been the less such because they were for the most part unaware that this was the role they were performing. In other words, they not only propagated but shared to the full in the illusion of universalism which was described earlier as the 'false consciousness' of the bourgeoisie, or of any class which cloaks its partial interests, deliberately or not, in the garb of universal principles, sacred and eternal verities, the national interest, and so on.

However, the Marxist view of intellectuals is rather more complex, even ambiguous, than this might suggest. To begin with, Marx and Engels, it will be recalled, said in the Communist Manifesto that, in 'times when the class struggle nears the decisive hour , a portion of the bourgeoisie goes over to the proletariat, and in particular a portion of the bourgeois ideologists'.\endnote{Revs., p. 77.} This notion was subsequently widened to include the possibility that some ideologists would join the socialist cause even in periods when the class struggle did not near its decisive hour; and this has more recently become a major element of the Marxist perspective in regard to the struggle against capitalism: with the vast increase in the numbers of people who perform 'intellectual' tasks under capitalism, and do so as 'pro- letarianized5 members of the 'collective labourer', the question of their alliance with the working class, and enlistment in its cause, becomes a matter of crucial importance. On the other hand, it is true that we are not here speaking of the 'bourgeois ideologists' that Marx and Engels had in mind, and are moving somewhat closer towards Gramsci's much wider definition.

As far as 'bourgeois ideologists' are concerned, Marx noted, as shown in the last quotation, that many (or some) of them might be particularly prone to 'go over to the proletariat' in times of crisis: they would do so because they had 'raised themselves to the level of comprehending theoretically the historical movement as a whole'.\endnote{Ibid., p. 77.} Marx and Engels also noted that these sections of the bourgeoisie that were 'proletarianized' would 'supply the proletariat with fresh elements of enlightenment and progress',\endnote{Ibid., p. 77.} and the point would presumably apply with particular force to 'ideologists', but this is no more than a presumption—Marx and Engels attributed no special 'role' to the intel- lectuals-cum-ideologists who had 'gone over' to the proletariat.

By the end of the nineteenth century, however, there were many such intellectuals occupying positions of influence in mass working-class parties, to the point of forming, so to speak, their own informal 'International' inside the Second International which had been formed in 1889. The question of intellectuals and their relation to the working-class movement had now assumed a much greater importance than it had ever had for Marx; and it was about them that Karl Kautsky wrote at the turn of the century what Lenin described in What is to be Done? as 'the following profoundly true and important words':
\begin{displayquote}
    \dots Socialism and the class struggle arise side by side and not one out of the other; each arises under different conditions. Modern socialist consciousness can arise only on the basis of profound scientific knowledge\dots The vehicle of science is not the proletariat but the bourgeois intelligentsia\dots it was in the minds of individual members of this stratum that modern socialism originated, and it was they who communicated it to the more intellectually developed proletarians who, in their turn, introduce it into the proletarian class struggle where conditions allow that to be done. Thus, socialist Consciousness is something introduced into the proletarian class struggle from without, and not something that arose within it spontaneously.\endnote{Lenin, What is to be Done?, CWL, vol. 5, p. 383. Italics in Kautsky's original text.}
\end{displayquote}

There is more than a faint echo of these 'profoundly true and important words' in Lenin's own insistence that 'the theory of Socialism . . . grew out of the philosophic, historical and economic theories elaborated by the educated representatives of the propertied classes, by intellectuals.'\footnote{Ibid., p. 375. Lenin also notes that 'in the very same way, in Russia, the theoretical doctrine of Social-Democracy arose altogether independently of the spontaneous growth of the working-class movement; it arose as a natural and inevitable outcome of the development of thought among the revolutionary socialist intelligentsia' (ibid., p. 375).}

Even so, the similarities with Kautsky are misleading. For Lenin had a poor view of intellectuals, not excluding Marxist ones. He was perfectly willing to acknowledge the debt which the development of socialist theory owed to bourgeois intellectuals. But he was exceedingly critical of what he saw as the intellectuals' 'flabbiness', individualism, lack of capacity for discipline and organization, tendency to reformism and opportunism, and so forth. He was also given to contrast their uncertain and vacillating commitment to the revolutionary movement with the sturdier and more reliable attitude of revolutionary workers. This is a 'workerist' tendency which has to a greater or lesser degree affected most Marxist writing and thinking. In Lenin's vision of the dictatorship of the proletariat, as expressed tor instance in The State and Revolution on the eve of the Bolshevik seizure of power, it is the working class which fills the stage, together with peasants and soldiers: 'experts' are expected to serve the revolution under the watchful (and distrustful) eye of the people. Nor was his perspective different after the Bolshevik seizure of power. In so far as intellectuals played a role in the leadership of the revolution, and a large number of them, in the Bolshevik party, did, they did so not as intellectuals but, precisely, as members of the party. Whatever Lenin thought of the merits of bourgeois intellectuals in an earlier epoch of socialist development, it is as members of the Party and through its mediation that he saw later generations of intellectuals involved in the service of the workers' cause.

This notion is further extended and articulated by Gramsci, for whom the emergence of a class of intellectuals 'belonging' to the working class was a matter of the first importance: in its bid for hegemony, the working class must produce its own intellectuals, and it is indeed one of the prime tasks of the revolutionary party to form such 'organic' intellectuals of the working class. These intellectuals play a major role in the formulation and the implantation of a new 'common sense'. But they play that role as members of a party which is itself the directing centre of the revolutionary movement, the 'modern prince', who is also the 'collective intellectual' of the working class. Here too it is the revolutionary party which performs the crucial task of bringing together and fusing in a common endeavour the working class and its intellectuals.

These and similar formulations do very little to resolve the real problems involved in the relation of intellectuals to the revolutionary movement, and this relation is in fact a major point of tension in Marxist theory, or at least in Marxist theory after Marx, and it has been a major point of tension, to put it mildly, in Communist practice.

Intellectuals are required to 'serve the people': this has been the injunction of Communist leaders to intellectuals from Lenin to the present day. The injuction may be entirely unobjectionable. The problems begin when the question is asked how the objective is to be served—not because the question is necessarily difficult, but because the answer to it is reserved, in the main, to the party leaders. It is they who have, in Communist practice, arrogated to themselves the prerogative of defining what serving the people means, and what it does not mean. This is much less true for Lenin, who tended to be cautious in such matters. But it has formed an intrinsic part of Communist practice ever since Lenin, and is in fact one characteristic feature of Stalinism, and for that matter of Maoism. Thus, one of the most authoritative of Chairman Mao's texts on the subject includes the following passage:
\begin{displayquote}
    Empty, dry dogmatic formulas do indeed destroy the creative mood; not only that, they first destroy Marxism. Dogmatic 'Marxism' is not Marxism, it is anti-Marxism. Then does not Marxism destroy the creative mood? Yes, it does. It definitely destroys creative moods that are feudal, bourgeois, petty-bourgeois, liberalistic, individualistic, nihilistic, art- for-art's sake, aristocratic, decadent or pessimistic, and every other creative mood that is alien to the masses of the people and to the proletariat. So far as proletarian writers and artists are concerned,  should not these kinds of creative moods be destroyed? I think they should; they should be utterly destroyed. And while they are being destroyed, something new can be constructed.\footnote{'Talks at the Yenan Forum on Literature and Art', in Selected Works of Mao Tse-tung (Peking, 1967), III, p. 94. Chairman Mao also noted that 'wrong styles of work still exist to a serious extent in our literary and art circles and that there are still many defects among our comrades, such as idealism, dogmatism, empty illusions, empty talk, contempt for practice and aloofness from the masses, all of which call for an effective and serious campaign of rectification' (ibid., p. 94). Such a comprehensive and vague indictment offers of course a formidable weapon against intellectuals, whose work is bound to be vulnerable to one or other such description.}
\end{displayquote}

This text dates from 1942, but there is plenty in Chairman Mao's pronouncements on intellectuals and intellectual and artistic work before and after that date which shows it to be typical of a mode of thought. The point about pronouncements such as these is not whether they are true, or false, or indeed whether they actually mean anything: it is rather that there is an external authority which is able to decide what is 'feudal, bourgeois, petty-bourgeois\dots{}' etc., and to impose its decisions upon intellectuals. The manner of imposition has varied, and assumed different forms in Russia under Stalin from those it has assumed in China, or in Cuba. But the prerogative of party leaders or of their representatives to decide what, in intellectual productions, 'serves the people', and even more so what does not, has not come under serious question in Communist regimes.

For a long time, that prerogative was also asserted, though it could not of course be imposed, by Communist Party leaders outside Communist regimes. The position has in recent years become much more open and fluid, and the notion of freedom of intellectual and artistic endeavour has by now gained a very considerable measure of acceptance and recognition by the leaderships, for instance, of the Italian and French Communist parties. The idea that intellectuals and artists should 'serve the people' has not been abandoned by these leaderships, as indeed why should it be? But the idea that they or some Cultural Commission of the Party should define what this does or does not involve has come to be increasingly discredited, together with the idea that the Party (meaning Party leaders) have any business to pass judgement on what, in artistic matters, is or is not 'acceptable'. This more relaxed approach to artistic and cultural work which has no direct and obvious political connotations has even come to affect, up to a point, intellectual work which is directly political and politically 'sensitive'. The general tendency is obviously towards a loosening up and the toleration of differences.

On the other hand, intellectuals in Western Communist parties are still expected, like other Party members, to behave in accordance with the principles of 'democratic centralism', at least where directly political utterances are involved. This may not render 'deviation' from party policy altogether unacceptable to party leaders, but it makes serious and sustained debate, and particularly organized debate, with defined platforms, next to impossible. This, which is one feature of the weakness or absence of internal democracy in these parties, affects and is intended to affect all Party members. But it affects Party intellectuals more than others in so far as criticism and probing are intrinsic elements of their being as intellectuals, or at least ought to be. The weight of 'democratic centralism' upon them is not as heavy as it was: but it is still very real; and the pressures and temptations to 'tighten up' which will arise when these parties come to office, and in the extremely difficult circumstances which will follow their accession to office, are bound to be very considerable.

Even so, the point may be worth making that the place and role of intellectuals is not a question determined by the nature of Marxism. It has been determined in Communist regimes—very adversely for the most part, namely at the price of conformity—by the conditions in which regimes which claim to be inspired by Marxism have come to power and developed. Ruthless repression of dissident intellectual work—and of every other kind of dissidence as well—has occurred in the name or 'under the banner' of Marxism. But there is no warrant for the view that such repression is somehow inscribed in the nature of Marxism, Assertions of this sort are merely propaganda.

This propaganda is part of the 'battle for consciousness' to which reference was made earlier, and which is part of the daily life of class societies. But this battle is only one facet of class conflict. There are other battles which proceed alongside this one, and in which the state plays a central role. Nor indeed is it really possible to separate any of these battles from each other; and the state intervenes, at one level or another, in all of them.

 Here is the first and the last defender of the old order; and it is therefore to the state that we must now turn.

 \theendnotes